% SPDX-License-Identifier: CC-BY-4.0
% Copyright (c) 2026 Martin Schröder <info@swedishembedded.com>

\section{Evaluation design}\label{sec:evaluation}

Three studies are pre-registered (\code{prereg/}): the teaching test, the in-conversation retraction test and the spoken conversations. The test split of the teaching study was registered before it was run; the retraction and spoken registrations were written after development runs of the same
scenarios, which they describe, and before the runs that count. The reading study, the extraction study, the training gate, the latency run and the recogniser screens are exploratory: their sets and checks were written before the readers or models were scored, but no registration file precedes them.
Sets are named by content digest. A held-out split is run once for a recipe; development runs are stated as such and are not results. Intervals are Wilson intervals over the asks (three per scenario, so correlated); the unit of resampling in the registration is the scenario, and at that level the intervals are wider.

\paragraph{Teaching scenarios (T2).}
Each scenario is a list of utterances that teach, optionally a ``forget that'', and asks with text checks written before any reply exists (\code{data/teaching}). The facts are about an invented company with values no model
can know, so a right answer is evidence of teaching and not of memory. Kinds: a fact, a corrected fact (taught, then corrected), a fact that must be used rather than recalled, sixteen facts at once, a procedure, a rule,
two rules at once, a fact taught and then forgotten, and a fact taught followed by unrelated questions. The same model answers twice: untaught, with an empty ledger, and taught, in a new session holding the ledger that a first
session built from the scenario's utterances, each confirmed with a yes. Twenty scenarios are dev and fifty are test. A second set asks the same questions in the \emph{same} session after ``forget that''.

\paragraph{Reading study (T1, T3).}
Authored utterances in many phrasings, a third in the unpunctuated lower case a recogniser produces, each with its gold kind (none, fact, rule, procedure) and a key word the content must hold; one in three is dev. A reading is right
for a teaching when its kind is the gold kind and its content holds the key, and right for a non-teaching when it is null.

\paragraph{Spoken conversations (T4).}
Scripted conversations whose user turns are rendered in a synthetic voice and recognised before the agent hears them, with pass conditions fixed before the run.

\paragraph{Training gate (S1).}
A reward with a known answer (fewer codec frames per word) checks that the optimiser moves a policy. The three methods share sentences, takes per update, seed and budget; the guard is the word error rate of the held-out takes under the streaming recogniser.

\paragraph{Intelligibility.}
Intelligibility is the word error rate under a recogniser. A one-shot transducer can cut the end of a sentence or return nothing, so a take is recognised with a tail of silence (\cref{sec:results}).
